\subsection{Relation to Q-ball excitations and continuum trapping}
\label{sec:related-work}

Q-balls belong to the established theory of localized scalar configurations
carrying a conserved global charge~\cite{Coleman1985}. Their perturbations
and nonlinear interactions have an extensive history. Battye and Sutcliffe
simulate collisions, charge transfer and fission in one, two and three
spatial dimensions in the same classical sextic model, after unit
conversion~\cite{BattyeSutcliffe2000}. Their kinetic normalization is
$\tfrac12|\partial\phi|^2$ and their potential is $2\mathcal U$;
$x=2x_{\mathrm{BS}}$, $t=2t_{\mathrm{BS}}$ give $\omega=\omega_{\mathrm{BS}}/2$.
Integrated energies and charges must also be converted in the chosen
spatial dimension. Their slowly decaying one-dimensional breathing
example is a dynamical precursor, not a certificate of an embedded linear
mode. The model itself is not new here. Its perturbations
require a coupled frequency problem. Kovtun, Nugaev, and Shkerin formulate
the complex-field linearization with conjugate coupling and the two
frequencies $\omega\pm\rho$~\cite{Kovtun2018}. Their three-dimensional
partial-wave equations provide a direct structural comparison, while the
vibrational modes selected in their piecewise-parabolic example satisfy
$\omega+\rho<1$. That calculation does not establish the embedded modes
of the sextic potential studied here.

Azatov, Ho, and Khalil study three-dimensional scattering in the same
sextic model family~\cite{Azatov2024}. With their $\phi_Q=\sqrt2 f$ and
$g=\beta=1/2$, their linear coefficients $U,W$ equal our $D,C$.
Conjugating their carrier convention and reducing the radial functions
maps their Appendix~B partial-wave equations to
Eqs.~\eqref{eq:radialA}--\eqref{eq:radialB}.
The radial operator is established prior work; the question here is the
existence of the particular localized solutions enclosed below.

The same work also treats Friedberg--Lee--Sirlin Q-balls containing a
complex and a real scalar. Its Section~3 derives the coupled charged
sidebands and neutral channel, their asymptotic thresholds, and scattering
relations~\cite{Azatov2024}. Thus an additional real field and a
three-channel linearization are not new ingredients. These scattering
calculations do not provide the particular infinite-domain embedded-mode
enclosures examined here.

Zhang, Zhou, and Zhu obtain analytic scattering solutions for stepwise
approximations to Q-balls of the same sextic
family~\cite{ZhangZhouZhu2025}. With their $g=\beta$ and
$f_0^2=1/(2\beta)$, their interior wavenumber $\sqrt{-\rho_1}$,
Eq.~(101), equals $k_{\mathrm{in}}$ in Eq.~\eqref{eq:interior}.
Their propagating scattering regime requires both exterior channels to
be open. The common interior dispersion and thin-wall phase structure
are therefore established ingredients; they do not themselves establish
cancellation of the remaining open-channel amplitude in the
one-open-channel regime considered here.

Ciurla, Dorey, Roma\'nczukiewicz, and Shnir study a sextic Q-ball model in
$1+1$ dimensions~\cite{Ciurla2024}. Their separation of bound,
half-propagating, and quasinormal modes is particularly relevant: in the
interval $1-\omega<\rho<1+\omega$, one asymptotic channel propagates and
the other is evanescent. They exhibit a small but nonzero radiation tail
associated with a long-lived quasinormal mode and investigate nonlinear
relaxation. This distinguishes a narrow resonance from the exact
vanishing of the propagating tail required here. Their one-dimensional
examples do not supply the spherical profile, origin conditions, or
infinite-domain enclosure used in the present three-dimensional problem.

Evslin et al. analyze paired Floquet modes of small-amplitude Q-balls in
$1+1$ dimensions~\cite{Evslin2026}. Their counterrotating analysis relates
a discrete component coupled to a propagating partner to a Feshbach-type
quasinormal mode. The small-amplitude approximation and one-dimensional
profile do not supply the full spherical sextic problem or the vanishing
open-channel amplitude examined here.

Continuum trapping through channel interference is also an established
mechanism. Friedrich and Wintgen relate interfering resonances to a
vanishing decay width~\cite{FriedrichWintgen1985}. A rigorous related
construction by Yu and Lu treats three Schr\"odinger channels on a
half-line, with one open channel, a crossing of two uncoupled bound
levels, and suitable nonzero radiation couplings~\cite{YuLu2025}.
Their small-coupling theorem concerns a potential that is exactly constant
outside a finite interval. It is not directly applicable to the
frequency-dependent two-channel Q-ball pencil, whose coefficients derive
from a nonlinear profile and retain an infinite tail. We therefore use
these works as comparisons, without assigning the specific
Friedrich--Wintgen mechanism to our modes solely from a vanishing
radiation amplitude.

Resonant total reflection through coupled propagating and evanescent
channels is established for discrete breathers~\cite{Flach2003},
optical solitons~\cite{Flach2005}, and spinor-condensate domain
walls~\cite{Watabe2012}. These are scattering solutions with an incident
wave; localization of a closed-channel component is not by itself an
embedded eigenmode of the full problem. Nor does the existence of open
and closed channels suffice for reflection: the fully coherent homogeneous
optical example remains transparent~\cite{Flach2005}.
These comparisons motivate a broader channel-interference interpretation,
without making an antiparticle branch a universal prerequisite or
establishing a necessary condition on Krein signatures.

Accumulation of embedded solitons also predates this work. Malomed et al.
derive a WKB/Bohr--Sommerfeld description in a quadratic optical model,
with a predicted scaling $\varepsilon_n\sim n^{-6/5}$~\cite{Malomed2005}.
Their matching is explicitly nonrigorous and its leading prefactor differs
from their numerical fit. Their objects are nonlinear embedded solitons;
ours are linear internal modes on changing Q-ball backgrounds. Neither
the scaling parameters nor their exponents can be identified merely from
the existence of a sequence of radiation cancellations.

Soffer and Weinstein prove nonlinear radiation damping under a positive
Fermi-golden-rule hypothesis for a real cubic Klein--Gordon
model~\cite{SofferWeinstein1999}. Its initial discrete mode lies below the
continuum and its third harmonic radiates. Their theorem does not directly
cover our embedded mode, rotating background, or second-harmonic source.
Our conditional continuation obstruction is not a dynamical decay theorem.

Suppressed radiation also occurs for real-field oscillons. Zhang et al.
report multiple decay-rate minima for the $\phi^6$ potential in
Section~7.5~\cite{ZhangOscillons2020}. Their analysis concerns long-lived
nonlinear configurations with radiative decay, rather than an enclosure
of an exact embedded eigenmode. The authors also identify limitations in
lifetime estimates near decay-rate minima. A small radiation signal or
a long-lived oscillation is therefore not, by itself, the criterion
used for the linear localization claims here.

An especially close comparison for the nonlinear radiation problem is
provided by Inagaki and Murakami~\cite{InagakiMurakami2026}.
They numerically resolve five zeros of the second-harmonic source overlap
for an internal mode of a thermal critical bubble, including three on
the thinner-wall side. The fundamental mode is below the continuum;
it is the leading second-harmonic radiation amplitude that vanishes.
At a selected zero, third-harmonic radiation remains and controls the
leading small-amplitude decay. These results neither constitute a linear
embedded eigenmode nor establish a radiation-free nonlinear solution.
They also do not establish an infinite sequence or its asymptotic spacing.
Our linear localization claim concerns the fundamental mode itself;
this distinction does not by itself establish priority.

The contribution examined here is the existence of particular localized
solutions of that full linearized Q-ball problem at the stated parameter
boxes, with regular origin conditions and controlled infinite tails.
The literature above does not replace those checks, and this comparison
does not establish priority. A finite collection of certified points also
does not imply an infinite ladder, or that modes on distinct backgrounds
belong to one fixed Q-ball spectrum. Finally, linear localization alone
does not establish an exactly localized nonlinear oscillation of finite
amplitude or nonlinear stability.
